\section{Related Work}
\label{sec:related}

\subsection{Measuring the quality of student test suites}
Many studies use coverage, mutants or faulty student submissions to grade test suites. Aaltonen et al.~\cite{aaltonen2010} found that most students reached the maximum coverage score while mutation scores were more spread out. Edwards and Shams~\cite{edwards2014} and Shams~\cite{shams2015} compared test quality measures using real faults in student programs as the reference and found that coverage is not strongly correlated with fault detection. Perretta et al.~\cite{perretta2022}, on 2,711 submissions, found that mutation score correlates strongly with detecting hand-seeded faults and fairly strongly with detecting faulty student submissions. Hall and Baniassad~\cite{hall2022} found that students graded by mutation analysis wrote more correct tests. These studies evaluate how test suites are graded. They do not compare the kind of feedback students see while writing tests.

\subsection{Feedback type as an intervention}
Mansur et al.~\cite{mansur2024} compared semesters that used only coverage with one semester that used mutation testing and found better test suites and code in the latter. Mansur et al.~\cite{mansur2026} went on to compare three semesters and reported that the original PIT report did not meet students' needs without modification and additional training. Alonge~\cite{alonge2025} randomly assigned 10 students to branch or mutation feedback and measured them with tests taken before and after, without feedback; the mutation group increased its rate of catching conceptual mutants by about 0.11 and the branch group by about 0.06. Shin and Kazerouni~\cite{shin2024} interviewed 12 students using think-aloud and found that mutation feedback easily causes cognitive overload. Cordova et al.~\cite{cordova2021} and Potter et al.~\cite{potter2025} showed that how feedback is presented (conceptual level, scaffolding, visualisation) affects outcomes. Overall, the evidence leans towards mutation feedback, but it comes from small samples, between-group or between-semester designs, and test suite quality as the outcome.

\subsection{Coverage, mutants and real faults}
Inozemtseva and Holmes~\cite{inozemtseva2014} found only a low to moderate correlation between coverage and test suite effectiveness when the number of tests is controlled for. Just et al.~\cite{just2014} found that detecting mutants is significantly correlated with detecting real faults, independently of coverage. Papadakis et al.~\cite{papadakis2018} found this correlation to be weak when test suite size is controlled for. For this reason, this study uses real faults as the outcome and includes test suite size in the analysis.

\subsection{Differences from the closest work}
Compared with Alonge~\cite{alonge2025}, this study asks a different question: how effectively tests are added while feedback is available, rather than skill transfer. It uses a within-person crossover design instead of two groups, uses Java and JUnit, and measures with real WA submissions instead of designer-made mutants. Shin and Kazerouni~\cite{shin2024} also had the same person use several kinds of feedback, but theirs is a qualitative study and does not measure detection of faulty submissions.
